\originalchapter{8.06 作业4：WKB、Airy与格点隧穿}
来源为Spring 2018 Assignment 4, 2018年3月1日发布, Barton Zwiebach. 第1--4、6题中文转述并给\textbf{AI独立解答}. 第3题保留loosely based on Griffiths署名, 按MIT完整任务解答; 第5题仅Griffiths 8.15题号与MIT连接提示, 无完整教材题面, 不补造. 通用Airy连接已在正文独立展开.
\originalsection{第1题：WKB到底解哪条方程}
\begin{exercise}
$-\hbar^2\psi''=p(x)^2\psi$, 近似$\psi_a=p^{-1/2}\exp[\ii\int^xp/\hbar]$. 求其精确满足的修正方程, 用de Broglie波长表达修正小的条件.
\end{exercise}
\begin{solution}
先对数微分$\psi_a'/\psi_a=-p'/(2p)+\ii p/\hbar$, 再平方加导数, 交叉虚项抵消, 得$\psi_a''/\psi_a=-p^2/\hbar^2-p''/(2p)+3p'^2/(4p^2)$. 因而
\[
-\hbar^2\psi_a''=\left[p^2+\frac{\hbar^2p''}{2p}-\frac{3\hbar^2p'^2}{4p^2}\right]\psi_a.
\]
各修正与$p^2$相比须小, 而不是彼此抵消就无条件相信近似. 令$\lambda=2\pi\hbar/p$, 两相对量为$|2\lambda'^2-\lambda\lambda''|/(8\pi^2)$和$3\lambda'^2/(16\pi^2)$. 常用足够的局部尺度条件$|\lambda'|\ll2\pi$, $|\lambda\lambda''|\ll(2\pi)^2$; 在$p=0$转折点它们失效, 要用Airy连接. 局部残差小仍不自动保证整个区间累积相位误差小.
\end{solution}
\originalsection{第2题：线性势的精确Airy谱}
\begin{exercise}
$V=gx$于$x>0$, 原点硬墙,$g>0$. (a)$L^3=\hbar^2/(2mg)$, 化为$\psi''(u)=u\psi$. (b)Fourier变换并给Ai积分, 核$\operatorname{Ai}(1)$. (c)求前两零点与能量.
\end{exercise}
\begin{solution}
令$\tilde u=x/L$, $\epsilon=E/(gL)$, 方程$\psi''=(\tilde u-\epsilon)\psi$, 再取$u=\tilde u-\epsilon$. Fourier约定$\tilde\psi(k)=\int\e^{-\ii ku}\psi(u)\dd u$使$-k^2\tilde\psi=\ii\partial_k\tilde\psi$, 所以$\tilde\psi=\e^{\ii k^3/3}$当$\tilde\psi(0)=1$. 逆变换将正负$k$配对给$\operatorname{Ai}(u)=\pi^{-1}\int_0^\infty\cos(k^3/3+ku)\dd k$, 振荡积分以合适截断/收敛因子理解. 数值$\operatorname{Ai}(1)=0.135292416313$, 与题给0.135292一致.

正$u$不增长的支是Ai, 舍Bi. 原点边界$\operatorname{Ai}(-\epsilon)=0$, 前两个负零点$-2.338107410460,-4.087949444131$. 所以$E_1=2.338107410460(\hbar^2g^2/(2m))^{1/3}$,$E_2=4.087949444131$乘同尺度. 单个硬壁+单个软转折点不应套两软转折点的$n+1/2$量子化.
\end{solution}
\originalsection{第3题：重力弹球的巨大量子数}
\begin{exercise}
势为$mgx$于$x>0$, $x=0$为硬墙. (a)用WKB求能量. (b)估计中子最低能量（eV）, 取$g=9.81$m/s$^2$. (c)1g物体从1m高度静止落下对应多大量子数? 原件署loosely based on Griffiths.
\end{exercise}
\begin{solution}
硬墙要求节点, 软转折点接到衰减支给$\pi/4$相位. 所以
\[
\int_0^{E/(mg)}\sqrt{2m(E-mgx)}\dd x
=\frac{2\sqrt{2m}E^{3/2}}{3mg}
=\pi\hbar(n+3/4),\qquad n=0,1,\ldots.
\]
解得$E_n=[3\pi\hbar g\sqrt m(n+3/4)/(2\sqrt2)]^{2/3}$. 中子质量约$1.67493\times10^{-27}$kg, 最低WKB为$1.39629\times10^{-12}$eV, 上一题Ai精确值为$1.40704\times10^{-12}$eV. 硬墙与软端不能一起用两软端的$n+1/2$.

经典能量$E=mgh=0.00981$J, 反解$n+3/4=2\sqrt{2m}E^{3/2}/(3mg\pi\hbar)$得$n\simeq8.91317\times10^{30}$. 这对应能量分布很窄的半经典波包的特征量子数; 对任意宽分布仅给平均能量不能反推精确平均$n$. 因而相邻能级相对变化$\Delta E/E\simeq2/(3n)$极小, 宏观能量近似连续; 不是认为量子规律在宏观物体上停止成立.
\end{solution}
\originalsection{第4题：四次势的半经典精度}
\begin{exercise}
$V=\alpha x^4$, $E_n=\gamma e_n$,$\gamma=(\hbar^4\alpha/m^2)^{1/3}$. 推转折点与两软端量子化, 得$\tilde e_n=\beta(n+\delta)^\epsilon$, 求$n=0$至5并与原数值$0.667986,2.39364,4.69680,7.33573,10.2443,13.3793$比较, 给0、2、5相对误差.
\end{exercise}
\begin{solution}
$x_0=(E/\alpha)^{1/4}$, 内动量$\sqrt{2m(E-\alpha x^4)}$. 两软端各$\pi/4$连接相位给$\int_{-x_0}^{x_0}p\dd x=(n+1/2)\pi\hbar$, 不含硬壁. 令$x=x_0t$, 积分$2\sqrt{2m}E^{3/4}\alpha^{-1/4}I$, $I=\int_0^1\sqrt{1-t^4}\dd t=\sqrt\pi\Gamma(1/4)/[8\Gamma(7/4)]\simeq0.874019184764$. 因而$\beta=[\pi/(2\sqrt2I)]^{4/3}$,$\delta=1/2,\epsilon=4/3$. 自有数值表见下, 低态相对误差较大, 越高能逐渐改善, 并不把WKB声称为精确谱.
\begin{center}\begin{tabular}{rrrr}\toprule $n$&WKB&原题数值&相对误差(\%)\\\midrule
0&0.54626733&0.667986&-18.221740\\
1&2.36356144&2.393640&-1.256603\\
2&4.67051993&4.696800&-0.559531\\
3&7.31480260&7.335730&-0.285280\\
4&10.22653643&10.244300&-0.173400\\
5&13.36376421&13.379300&-0.116118\\
\bottomrule\end{tabular}\end{center}

\end{solution}
\originalsection{第6题：高阶微扰怎样产生指数隧穿}
\begin{exercise}
格点$0,\ldots,N$两端势0、内部$V_0$, 最近邻跃迁$-\lambda$, $\lambda\ll V_0$. (a)原两端零能态何阶第一次分裂, 求间隙. (b)由步长$\ell$二阶差分动能求$\lambda$; 固定$\lambda/V_0=\varepsilon\ll1$, 验证$\ell\to0$时$\ell^2V_0$保持不变. 再用$\ell=W/N$消去$\ell$, 求间隙随$\varepsilon,W,V_0,m,\hbar$的尺度. (c)只保两态, 初始$\ket0$首次完全到$\ket N$何时, 转移频率怎样与WKB比较?
\end{exercise}
\begin{solution}
连两端最短路径要$N$跳, 更低阶只有相同对角位移, 不分裂. 唯一路径非对角有效元$t=(-\lambda)^N/(-V_0)^{N-1}=-\lambda^N/V_0^{N-1}$. 偶组合$g_+$低, 奇组合$g_-$高, 间隙$\Delta=2\lambda^N/V_0^{N-1}$至首个非零阶, 包括$N=1$直接耦合特例. 回程修正为更高阶.

差分$-\hbar^2\partial_x^2/(2m)$的非对角系数$-\hbar^2/(2m\ell^2)$, 所以$\lambda=\hbar^2/(2m\ell^2)$. 忽略的对角项是共同常能位移, 不改变分裂, 但固定物理$V_0$时$\ell\to0$会使$\lambda/V_0$增大, 微扰失效. 设$\varepsilon$固定意味着$\ell^2V_0=\hbar^2/(2m\varepsilon)$固定, 这是同步变高障壁的尺度族, 不是固定势的连续极限. 消去给
\[
\Delta\simeq2V_0\exp\!\left[-\frac{W\sqrt{2mV_0}}\hbar\sqrt\varepsilon\log(1/\varepsilon)\right].
\]
两态演化从$\ket0=(g_++g_-)/\sqrt2$出发,$P_N(t)=\sin^2(\Delta t/(2\hbar))$. 首次完全转移$t_* =\pi\hbar/\Delta$, 若以$1/t_*$称“隧穿率”, 为$\Delta/(\pi\hbar)$且同指数. 这是相干首次转移频率, 概率不是恒定率的线性增长. WKB\textbf{通量透射概率}$\exp[-2W\sqrt{2mV_0}/\hbar]$同样随$W\sqrt{V_0}$指数抑制, 但比较的是不同量且此固定$\varepsilon$近似系数为$\sqrt\varepsilon\log(1/\varepsilon)$, 不能把两数等同为精确连续结果.
\end{solution}
